\section*{Pre-Registered Hypotheses and Outcomes}

\begin{table}[htbp]
\centering
\caption{Every pre-registered hypothesis and its verdict, with what each verdict
does and does not mean. The analysis plan was fixed before any experiment was run.
Several ``supported'' verdicts are weak, and the last column says why rather than
leaving it to the reader.}
\label{tab:hypotheses}
\small
\fittable{\begin{tabular}{lp{3.4cm}p{2.9cm}rrlp{3.6cm}}
\toprule
ID & Hypothesis & Test & Result & $p_{\text{adj}}$ & Verdict & Qualification \\
\midrule
H1 & Reliability-gated correction improves accuracy over a non-correcting agent. & series-mean change vs never (\%, positive = worse); Wilcoxon Holm-adj. & 0.421 & 0.005 & not supported & the effect is a small, significant degradation \\
H2 & Gains concentrate under distribution shift, anomalies and regime change. & difference-in-differences over 48 perturbed strata (Holm) & 2.000 & 0.005 & supported & 2 of 48 strata after Holm; both unconditional correction, mildest severity \\
H3 & Forecast failure is predictable from pre-hoc diagnostics. & AUROC of the risk estimator & 0.793 & --- & supported & supports triage, not correction; the gate built on it does not win \\
H4 & Selective tool use matches exhaustive use at lower cost. & mean action evaluations per instance & 0.088 & --- & supported & because the gate rarely acts, not because it acts well \\
H5 & Reliability-gated confidence is better calibrated than an agent's verbalised confidence. & paired bootstrap over series on the ECE difference & 0.220 & --- & supported & against the model's verbalised confidence, an easy comparator \\
H6 & Gains are not explained by extra compute alone (vs. rate-matched random correction). & RGC vs. random\_correct & 0.023 & --- & supported & because the gate rarely acts; it does not beat never correcting \\
H7 & A quantitative gate decides better than verbal self-critique (pre-registered comparison). & mean oracle regret, calibrated gate vs. LLM critique gate & 0.094 & --- & supported & the gate's regret is close to never correcting's; it wins by not acting \\
\bottomrule
\end{tabular}
}
\end{table}

\section*{Broader Impact Statement}

The practical message of this work is cautionary, and we think that is where its impact
lies. Automated self-correction is being added to forecasting systems on the assumption
that a model inspecting its own output will improve it. We find, in a controlled setting
with counterfactual ground truth, that unconditional correction reliably \emph{degrades}
point accuracy, and that the adaptive machinery intended to prevent that does not
establish an improvement over doing nothing. Forecasts drive decisions about inventory,
staffing, capacity and public health resourcing; a correction stage that silently makes
a subset of them worse, while presenting as a reliability feature, is a realistic way for
harm to enter such a system unnoticed. We would rather this result be available to
practitioners than not.

The positive finding carries its own risk. We show that a calibrated failure-risk
estimate is a useful triage signal for routing forecasts to human review. Triage systems
can entrench the disparities in the data they are fitted to: series that are short,
irregular, or from under-represented sources may be systematically over- or
under-flagged, and a reviewer working from a ranked queue has little visibility into
that. The risk estimator here is fitted on statistical diagnostics with no demographic
content, but the series themselves are not a neutral sample, and we have not audited
flagging rates across data sources. Anyone deploying such a signal should.

We see no dual-use concern specific to this work. The datasets are public, contain no
personal data, and are used under the licences recorded in the released manifest. The
release is intended to make correction policies, including ones that disagree with our
conclusions, cheap to evaluate against the same counterfactual outcomes.

\section*{Reproducibility Statement}

Every result is produced by a single command against a versioned configuration file, and
each run writes a manifest recording the configuration, its hash, the code version, the
environment and package versions. Reruns with the same configuration and seed reproduce
identical metrics. LLM responses are cached under a hash of model, prompt and parameters,
so the reported results can be regenerated without an API key and without network
access. All figures and tables in this manuscript are generated from results files by
\texttt{scripts/make\_paper\_assets.py}; no experimental number is entered by hand. The
pre-registered analysis plan, including all deviations from it, is in
\texttt{docs/ANALYSIS\_PLAN.md} in the released code.

\section*{Data Availability}

The benchmark generators, the counterfactual outcome tensor, the evaluation framework
and all analysis code are provided as an anonymised archive with this submission, and
will be released publicly under a permissive licence with an archived DOI on acceptance.
The archive includes the cached language-model responses, so every LLM arm in this paper
replays without an API key and without a provider account.

Real datasets are not redistributed. They are obtained from their original sources by
\texttt{scripts/prepare\_data.py}, which records the URL, licence, SHA-256 checksum and
access date of every file in \texttt{datasets/LICENSES.md}. Monash archive datasets are
CC~BY~4.0 \citep{monash2021}; the ETT data is used under the licence stated in its
upstream repository \citep{informer2021}, and only derived metrics are reported.

Synthetic data is not distributed as files because it is fully determined by its
generator, seed and index, and is regenerated deterministically by the released code.

\section*{Author Contributions}

\ifcameraready
Sreevallabh Kakarala is the sole author and carried out every role:
conceptualisation, methodology, software, validation, formal analysis,
investigation, data curation, visualisation, writing of the original draft, and
review and editing.
\else
This paper has a single author, who carried out every role: conceptualisation,
methodology, software, validation, formal analysis, investigation, data curation,
visualisation, writing of the original draft, and review and editing. The name is
withheld for anonymous review.
\fi
